% This file is generated from the corresponding Obsidian note.
% Source: Teoricos/GS - Teo21.md
\chapter{Formas diferenciales}
\label{ch:teo-21}
\begin{bookdefinition}{Fibrado de $k$-formas alternantes}
Sea $M$ variedad suave. Se define el fibrado de $k$-formas alternantes, denotado por $\Lambda^kT^*M$, como la union disjunta de todos los espacios de $k$-tensores alternantes sobre cada $T_p^{*}M$ con $p\in M$:

\[
\Lambda^kT^*M:=\bigsqcup_{p\in M}\Lambda^k(T_p^*M)=\{(p,\omega):p\in M\text{ y }\omega\in\Lambda^k(T_p^*M)\}.
\]
El caso $k=1$, es

\[
\Lambda T^*M=\{(p,\alpha):p\in M\text{ y }\alpha\in T_p^*M\}
\]
es llamado fibrado cotangente. Tambien, tenemos una proyeccion canonica

\[
\pi:\Lambda^k(T^*M)\to M,\qquad (p,\omega)\mapsto p.
\]
\end{bookdefinition}
\begin{bookremark}{}
Como sucedió con el fibrado tangente, es similar ver que $\Lambda^kT^*M$ es una variedad suave de dimension $\dim M+\binom{\dim M}{k}$.

\end{bookremark}
\begin{bookproof}{}
\begin{enumerate}
\item La idea es imitar lo que se hizo con el fibrado tangente. Si $n=\dim M$, para una carta $(U,\varphi=(x_1,\ldots,x_n))$, la carta inducida es
\end{enumerate}
\[
(\widetilde U,\widetilde\varphi=(x_1,\ldots,x_n,(v_{i_1\ldots i_k})_{1\le i_1<\cdots<i_k\le n})),
\]
donde $\widetilde U=\pi^{-1}(U)$ y, si $(p,\omega)\in\widetilde U$,

\[
v_{i_1\ldots i_k}(\omega)=\omega\left(\frac{\partial}{\partial x_{i_1}}\bigg|_p,\ldots,\frac{\partial}{\partial x_{i_k}}\bigg|_p\right).
\]
veamos por que esto ultimo vale. 2. Como $\left\{\frac{\partial}{\partial x_1}\big|_p,\ldots,\frac{\partial}{\partial x_n}\big|_p\right\}$ es base de $T_pM$, sabemos que $\{dx_1|_p,\ldots,dx_n|_p\}$ es base de $T_p^*M$. 3. Recordar que, formalmente, $dx_k|_p:T_pM\to T_{x_k(p)}\mathbb R$ y, por conveniencia, identificamos $(dx_k)_p$ con un escalar, y asi esta justificado que $(dx_k)_p\in T_p^*M$. 4. Además, por dualidad, se tiene

\[
dx_j|_p\left(\frac{\partial}{\partial x_i}\bigg|_p\right)=\delta_{ij}.
\]
\begin{enumerate}[start=5]
\item Y vimos que
\end{enumerate}
\[
\{dx_{i_1}|_p\wedge\cdots\wedge dx_{i_k}|_p\}_{1\le i_1<\cdots<i_k\le n}
\]
es base de $\Lambda^k(T_p^*M)$. Por tanto

\[
\omega=\sum_{1\le i_1<\cdots<i_k\le n}C_{i_1\ldots i_k}\,dx_{i_1}|_p\wedge\cdots\wedge dx_{i_k}|_p.
\]
\begin{enumerate}[start=6]
\item Para averiguar el coeficiente $C_{i_1\ldots i_k}$, evaluamos en la $k$-upla asociada $\left(\frac{\partial}{\partial x_{i_1}}\big|_p,\ldots,\frac{\partial}{\partial x_{i_k}}\big|_p\right)$.
\item En efecto, si $1\le j_1<\cdots<j_k\le n$, entonces
\end{enumerate}
\[
(dx_{j_1}|_p\wedge\cdots\wedge dx_{j_k}|_p)\left(\frac{\partial}{\partial x_{i_1}}\bigg|_p,\ldots,\frac{\partial}{\partial x_{i_k}}\bigg|_p\right)=\det\left(dx_{j_a}|_p\left(\frac{\partial}{\partial x_{i_b}}\bigg|_p\right)\right)_{a,b=1}^k.
\]
\begin{enumerate}[start=8]
\item Usando $dx_{j_a}|_p\left(\frac{\partial}{\partial x_{i_b}}\big|_p\right)=\delta_{j_a i_b}$, queda
\end{enumerate}
\[
(dx_{j_1}|_p\wedge\cdots\wedge dx_{j_k}|_p)\left(\frac{\partial}{\partial x_{i_1}}\bigg|_p,\ldots,\frac{\partial}{\partial x_{i_k}}\bigg|_p\right)=\det(\delta_{j_a i_b}).
\]
\begin{enumerate}[start=9]
\item Como los indices estan ordenados de forma creciente, este determinante vale $1$ si $(j_1,\ldots,j_k)=(i_1,\ldots,i_k)$, y vale $0$ en caso contrario.
\item Por lo tanto, al evaluar la expansion de $\omega$ en $\left(\frac{\partial}{\partial x_{i_1}}\big|_p,\ldots,\frac{\partial}{\partial x_{i_k}}\big|_p\right)$, todos los terminos se anulan excepto el termino correspondiente a $(i_1,\ldots,i_k)$, y obtenemos
\end{enumerate}
\[
\omega\left(\frac{\partial}{\partial x_{i_1}}\bigg|_p,\ldots,\frac{\partial}{\partial x_{i_k}}\bigg|_p\right)=C_{i_1\ldots i_k}.
\]
\begin{enumerate}[start=11]
\item Entonces
\end{enumerate}
\[
v_{i_1\ldots i_k}(\omega):=\omega\left(\frac{\partial}{\partial x_{i_1}}\bigg|_p,\ldots,\frac{\partial}{\partial x_{i_k}}\bigg|_p\right)
\]
es exactamente la coordenada de $\omega$ con respecto al elemento de base $dx_{i_1}|_p\wedge\cdots\wedge dx_{i_k}|_p$. 12. O sea que las ultimas $\binom nk$-funciones de $\widetilde\varphi$ nos dan las coordenadas de $\omega$ con respecto a la base

\[
\{dx_{i_1}|_p\wedge\cdots\wedge dx_{i_k}|_p\}_{1\le i_1<\cdots<i_k\le n}.
\]
\begin{enumerate}[start=13]
\item En esta estructura de variedad diferenciable, la proyeccion canonica $\pi:\Lambda^k(T^*M)\to M$ es suave.
\end{enumerate}
\bookqed
\end{bookproof}
\begin{bookdefinition}{k-forma alternante o simplemente k-forma}
Una $k$-forma $\omega$ es cualquier seccion de $\pi:\Lambda^k(T^*M)\to M$ es decir, es una funcion

\[
\omega:M\to\Lambda^k(T^*M)
\]
tal que $\pi(\omega(p))=p$, o lo que es lo mismo $\omega(p)=(p,\omega_p)$ con $\omega_p\in\Lambda^k(T_p^*M)$. Como hicimos con los campos, identificamos $\omega(p)$ con su segunda componente $\omega_p$.

\end{bookdefinition}
\begin{bookdefinition}{k-forma suave}
Decimos que una $k$-forma es suave cuando la sección de $\pi$,

\[
\omega:M\to\Lambda^k(T^*M)
\]
es una función suave. Se suele denotar el conjunto de tales secciones suaves por

\[
\Omega^k(M)
\]
o, en notación de fibrados vectoriales, por $\Gamma(\Lambda^k(T^*M))$.

\end{bookdefinition}
\begin{bookremark}{Para memorizar (intuicion)}
$\Omega^{k}(M)$ vendria a ser como el equivalente a $\mathfrak{X}(M)$ para campos Y cuando evaluas un $w\in \Omega^{k}(M)$ en un punto $p$ ahi obtenes $w_{p}$ un $k$-tensor alternante asi como cuando evaluavas un $X\in \mathfrak{X}(M)$ en $p$ obtenias $X_{p}$ un vector tangente (en ambos casos estoy usando la identificacion)

\end{bookremark}
\begin{bookremark}{}
Recordar que $\Lambda^0(V^*)=\mathbb R$ \hyperref[blk:2c0c65]{Teórico 20} por lo tanto

\[
\Lambda^0(T^*M)=\mathbb \bigsqcup_{p\in M}\{ (p,r):p\in M,\ r\in \Lambda^{0}(T_{p}^{*} M)=\mathbb{R}  \}=M\times \mathbb{R}
\]
Luego $\Omega^0(M)$ se identifica con $C^\infty(M)$. De la siguiente manera Si $\omega\in\Omega^0(M)$, entonces $\omega:M\to\Lambda^0(T^*M)$, $\omega(p)=(p,r_{p})$ con $r_{p}\in\mathbb R$, y $\omega$ es suave sii

\[
f:M\to\mathbb R \qquad p\mapsto r_{p}
\]
es suave. Pues aca $\Lambda^0(T^*M)=M\times\mathbb R$, y como podemos identificar $w(p)=w_{p}=r_{p}$ entonces tendriamos que $w=f$ como $w$ es suave $f$ es suave, osea $w=f\in C^{\infty}(M)$

\end{bookremark}
\section{Conjunto de formas diferenciales}
\begin{bookremark}{}
Se define el conjunto $\Omega^*(M)$ de las formas diferenciales, sin el $k$, sobre $M$ como la suma directa

\[
\Omega^*(M):=\bigoplus_{k=0}^n\Omega^k(M).
\]
Veremos que podemos llevar el producto cuña a este conjunto y dotarlo de estructura de algebra graduada.

\end{bookremark}
\section{Suavidad de $k$-formas}
\begin{bookremark}{}
Como nos sucedio en el capitulo de Campos, la primera pregunta es: criterios para saber cuando una $k$-forma es suave.

\end{bookremark}
\begin{bookproposition}{Coordenadas de $k$-formas}
Tenemos el resultado analogo de Campos. Sea $(U,\varphi=(x_1,\ldots,x_n))$ una carta suave. Para todo $p\in U$ tenemos que:

\begin{itemize}
\item Base de $T_pM$: $B=\left\{\frac{\partial}{\partial x_1}\big|_p,\ldots,\frac{\partial}{\partial x_n}\big|_p\right\}$.
\item Base de $T_p^*M$: $B^*=\{dx_1|_p,\ldots,dx_n|_p\}$.
\item Base de $\Lambda^k(T_p^*M)$: $B_k=\{dx_{i_1}|_p\wedge\cdots\wedge dx_{i_k}|_p:1\le i_1<\cdots<i_k\le n\}$.
\end{itemize}
Por tanto, si $\omega$ es una $k$-forma tenemos

\[
\omega_p=\sum_{1\le i_1<\cdots<i_k\le n}C_{i_1\ldots i_k}(p)\,dx_{i_1}|_p\wedge\cdots\wedge dx_{i_k}|_p,
\]
lo cual define funciones

\[
C_{i_1\ldots i_k}:U\to\mathbb R \qquad p\mapsto C_{i_1\ldots i_k}(p)
\]
que se llaman las coordenadas de $\omega$ con respecto al marco coordenado $\{dx_{i_1}\wedge\cdots\wedge dx_{i_k}\}$.

\end{bookproposition}
\phantomsection\label{blk:70585b}
\begin{bookremark}{}
Ademas, tenemos la siguiente construccion usando campos suaves. Sea $\omega$ una $k$-forma, $k\ge1$, y sean $X_1,\ldots,X_k$ campos suaves. Definimos

\[
\omega(X_1,\ldots,X_k):M\to\mathbb R,\qquad p\mapsto\omega_p((X_1)_p,\ldots,(X_k)_p).
\]
\end{bookremark}
\begin{bookremark}{}
Con estas consideraciones, muy similar a las estudiadas en campos, tenemos el siguiente resultado analogo cuya prueba es similar.

\end{bookremark}
\begin{booktheorem}{Suavidad de $k$-formas}
Sea $\omega:M\to\Lambda^k(T^*M)$ una $k$-forma. Las siguientes afirmaciones son equivalentes:

\begin{itemize}
\item (a) $\omega$ es suave.
\item (b) Para toda carta $(U,\varphi=(x_1,\ldots,x_n))$, las coordenadas de $\omega$ con respecto al marco coordenado $\{dx_{i_1}\wedge\cdots\wedge dx_{i_k}\}_{1\le i_1<\cdots<i_k\le n}$ son funciones suaves.
\item (c) Para cualquier $k$ campos suaves $X_1,\ldots,X_k\in\mathfrak X(M)$, la funcion $\omega(X_1,\ldots,X_k)\in C^\infty(M)$.
\end{itemize}
\end{booktheorem}
\begin{bookproof}{}
\begin{itemize}
\item $(a\Rightarrow b)$.
\end{itemize}
\begin{enumerate}
\item Supongamos que $\omega:M\to\Lambda^k(T^*M)$ es suave.
\item Sea $(U,\varphi=(x^1,\ldots,x^n))$ una carta de $M$.
\item Esta carta induce una carta $(\widetilde U,\widetilde\varphi)$ de $\Lambda^k(T^*M)$, donde $\widetilde U=\pi^{-1}(U)$. Usamos la notación de multiíndices $I=(i_1,\ldots,i_k)$, con $i_1<\cdots<i_k$, y escribimos $dx^I=dx^{i_1}\wedge\cdots\wedge dx^{i_k}$. La carta inducida está dada por
\end{enumerate}
\[
\widetilde\varphi\left(\sum_I C_I\,dx^I|_p\right)=\left(x^1(p),\ldots,x^n(p),(C_I)_I\right).
\]
\begin{enumerate}[start=4]
\item Como $\omega$ es suave, su restricción $\omega|_U:U\to\widetilde U$ es suave. Por lo tanto, como $\widetilde\varphi$ es suave, la composición
\end{enumerate}
\[
\widetilde\varphi\circ\omega|_U:U\to\mathbb R^{n+\binom nk}
\]
es suave. 5. Escribimos, para cada $p\in U$,

\[
\omega_p=\sum_I C_I(p)\,dx^I|_p.
\]
\begin{enumerate}[start=6]
\item Entonces
\end{enumerate}
\[
(\widetilde\varphi\circ\omega|_U)(p)=\left(x^1(p),\ldots,x^n(p),(C_I(p))_I\right).
\]
\begin{enumerate}[start=7]
\item Como $\widetilde\varphi\circ\omega|_U$ es suave, todas sus funciones coordenadas son suaves. En particular, para cada multiíndice $I=(i_1,\ldots,i_k)$ con $1\le i_1<\cdots<i_k\le n$, la función $C_I:U\to\mathbb R$ es suave.
\item Como la carta $(U,\varphi)$ era arbitraria, las coordenadas $C_I$ de $\omega$ respecto de todo marco coordenado $\{dx^I\}_I$ son funciones suaves.
\end{enumerate}
\begin{itemize}
\item $(b\Rightarrow a)$.
\end{itemize}
\begin{enumerate}
\item Supongamos que, para toda carta $(U,\varphi=(x^1,\ldots,x^n))$, las coordenadas $C_I:U\to\mathbb R$ de $\omega$ son funciones suaves.
\item Fijemos una de estas cartas y consideremos la carta inducida $(\widetilde U,\widetilde\varphi)$ de $\Lambda^k(T^*M)$, con $\widetilde U=\pi^{-1}(U)$.
\item Para cada $p\in U$, escribimos
\end{enumerate}
\[
\omega_p=\sum_I C_I(p)\,dx^I|_p.
\]
\begin{enumerate}[start=4]
\item Por la definición de la carta inducida,
\end{enumerate}
\[
(\widetilde\varphi\circ\omega|_U)(p)=\left(x^1(p),\ldots,x^n(p),(C_I(p))_I\right).
\]
\begin{enumerate}[start=5]
\item Las funciones $x^1,\ldots,x^n$ son suaves y, por hipótesis, todas las funciones $C_I$ son suaves. Por lo tanto, $\widetilde\varphi\circ\omega|_U$ es suave.
\item Como $\widetilde\varphi$ es una carta, concluimos que $\omega|_U:U\to\widetilde U$ es suave.
\item Como esto vale para toda carta $(U,\varphi)$ de $M$, concluimos que $\omega:M\to\Lambda^k(T^*M)$ es suave.
\end{enumerate}
\begin{itemize}
\item $(b\Rightarrow c)$
\end{itemize}
\begin{enumerate}
\item Tomamos un $p\in M$ ahora tenemos $(U,\phi=x_{1},\ldots,x_{n})$ carta entonces por (b)
\end{enumerate}
\[
\omega|_U=\sum_{i_1<\cdots<i_k}C_{i_1\cdots i_k}\,dx_{i_1}\wedge\cdots\wedge dx_{i_k},\qquad C_{i_1\cdots i_k}\in C^\infty(U)
\]
\begin{enumerate}[start=2]
\item Ademas
\end{enumerate}
\[
X_{i}=\sum^{n}_{j=1}a_{ij}\frac{\partial}{\partial x_{j}}
\]
\begin{enumerate}[start=3]
\item Luego
\end{enumerate}
\[
\omega|_{U}(X_{1},\ldots X_{k})=\sum_{i_1<\cdots<i_k}C_{i_1\cdots i_k}\,dx_{i_1}\wedge\cdots\wedge dx_{i_k}\left(\sum^{n}_{j=1}a_{1j}\frac{\partial}{\partial x_{j}},\ldots,\sum^{n}_{j=1}a_{kj}\frac{\partial}{\partial x_{j}}\right)
\]
\begin{enumerate}[start=4]
\item Ahora notamos que
\end{enumerate}
\[
\begin{aligned} dx_{i_1}\wedge\cdots\wedge dx_{i_k}\left(\sum^{n}_{j=1}a_{1j}\frac{\partial}{\partial x_{j}},\ldots,\sum^{n}_{j=1}a_{kj}\frac{\partial}{\partial x_{j}}\right)&=\det\left( dx_{i_{r}}\left( \sum^{n}_{j=1}a_{sj}\frac{\partial}{\partial x_{j}} \right) \right)_{r,s=1}^{k} \\ & =\det\left( \sum^{n}_{j=1} a_{sj}dx_{i_{r}}\left(\frac{\partial}{\partial x_{j}}\right) \right)_{r,s=1}^{k}\\&=\det(a_{s i_{r}})_{r,s=1}^{k} \end{aligned}
\]
\begin{enumerate}[start=5]
\item Pero entonces
\end{enumerate}
\[
\omega|_{U}(X_{1},\ldots,X_{k})(p)=(\omega|_{U})_{p}(X_{1},\ldots,X_{k})=\sum_{i_{1}<\ldots<i_{k}} C_{i_{1}\ldots i_{k}}(p)\det(a_{s i_{r}}(p))
\]
\begin{enumerate}[start=6]
\item Pero $C_{i_{1}\ldots i_{k}}$ es suave como funcion de $p$ por hipotesis, lo mismo para $\det(a_{s i_{r}})_{r,s=1}^{k}$ como funcion de $p$ por que $X_{i}$ son campos suaves. Y $\det$ es un polinomio de sus entradas.
\item Entonces $\omega|_{U}(X_{1},\ldots,X_{k})(p)$ es suave y obviamente como esto lo podemos hacer para cualquier $p$ tenemosq ue
\end{enumerate}
\[
\omega(X_{1},\ldots,X_{k})(p)
\]
es suave

\begin{itemize}
\item $(c\Rightarrow b)$
\end{itemize}
\begin{enumerate}
\item Tomamos una carta $(U,x_{1},\ldots,x_{n})$ como $B_k=\{dx_{i_1}|_p\wedge\cdots\wedge dx_{i_k}|_p:1\le i_1<\cdots<i_k\le n\}$ es base de $\Lambda^{k}(T_{p}^{*}M)$ tenemos
\end{enumerate}
\[
\omega|_{U}=\sum_{i_{1}<\ldots<i_{k}} C_{i_{1}\ldots i_{k}}\ .dx_{i_{1}}\wedge\ldots\wedge dx_{i_{k}}
\]
con $C_{i_{1}\ldots i_{k}}:U \rightarrow\mathbb{R}$. Bastaria mostrar que son suaves 2. Ahora tenemos que

\[
\begin{aligned}\omega|_{U}\left(\frac{\partial}{\partial x_{j_{1}}},\ldots, \frac{\partial}{\partial x_{j_{k}}}\right)(p)& =\omega_{p}\left(\frac{\partial}{\partial x_{j_{1}}}\bigg|_{p},\ldots, \frac{\partial}{\partial x_{j_{k}}}\bigg|_{p}\right)\\&=\sum^{\infty}_{i_{1}<\ldots<i_{k}}C_{i_{1}\ldots i_{k}}(p) \ .dx_{i_{1}}\wedge\ldots\wedge dx_{i_{k}}\left(\frac{\partial}{\partial x_{j_{1}}}\bigg|_{p},\ldots, \frac{\partial}{\partial x_{j_{k}}}\bigg|_{p}\right)\\&=C_{I}(p)\end{aligned}
\]
con $I=J$ osea el unico caso donde el determinante no es $0$ que es cuando coinciden los multi indices, $j_{1}=i_{1}\ldots j_{k}=i_{k}$ 3. Ahora lo unico que tenemos que corregir es que $\frac{\partial}{\partial x_{i}}$ son campos suaves sobre $U$  y para usar la hipotesis las necesito suaves sobre $M$ 4. Tomamos $q\in U$ y ahora tenemos $V_{q}\subseteq U$ abierto de $q$ y $f\in C^{\infty}(M)$ tal que $f|_{\overline{V_{q}}}\equiv 1$ y $\operatorname{supp}f\subseteq U$ y defino

\[
X_{i}(p)=\begin{cases} f(p)\frac{\partial}{\partial x_{i}}\bigg|_{p} & p\in U \\0 & x\not\in U \end{cases}
\]
que sabemos es un campo suave sobre todo $M$ y cerca de $p$ (osea en algun abierto dentro de $\overline{V_{q}}$, por ejemplo $V_{q}$ ) vale exactamente $\frac{\partial}{\partial x_{i}}\bigg|_{q}$ 5. Luego haciendo la misma cuenta que en 2. pero con $X_{i}$ tenemos que

\[
\sum_{i_{1}<\ldots<i_{k}} C_{i_{1}\ldots i_{k}}(q)
\]
es suave para todo $q\in V_{p}$ 6. Y como esto lo podemos hacer para cualquier $q\in U$ es suave en todo $U$ como queriamos

\bookqed
\end{bookproof}
\begin{bookproposition}{Suavidad con una sola carta (analogo a campos)}
Sea $(U,\varphi=(x_1,\ldots,x_n))$ una carta suave y sea $\omega$ una $k$-forma sobre $U$,

\[
\omega:U\to\Lambda^k(T^*M)
\]
entonces $\omega$ es suave si y solo si las coordenadas de $\omega$ con respecto al marco $\{dx_{i_1}\wedge\cdots\wedge dx_{i_k}\}_{1\le i_1<\cdots<i_k\le n}$ son funciones suaves.

\end{bookproposition}
\begin{bookproof}{}
3. Dada la carta $(U,\varphi)$ tenemos la carta inducida $(\widetilde U,\widetilde\varphi)$ de $\Lambda^k(T^*M)$ que es un difeomorfismo entre $\widetilde U$ y $U\times\mathbb R^{\binom nk}$

\[
\widetilde\varphi:\widetilde U\subset\Lambda^k(T^*M)\to \varphi(U)\times\mathbb R^{\binom nk}\subseteq \mathbb{R}^{n} \times\mathbb R^{\binom nk}
\]
\begin{enumerate}[start=4]
\item Y justo $\widetilde\varphi$ se encarga de dar las coordenadas de $\omega$ en el marco $\{dx_{i_1}\wedge\cdots\wedge dx_{i_k}\}$. Por tanto, la composicion
\end{enumerate}
\[
\widetilde\varphi\circ\omega: U\rightarrow \varphi(U)\times\mathbb R^{\binom nk}\qquad p\mapsto(\varphi(p),C_{i_{1}\ldots i_{k}}(p))
\]
(Notar que abuse notacion por que deberian ser todas las coordenadas, osea los subindices $i_{k}$ deberian moverse por todas las posibilidades, siempre ordenados de menor a mayor) 5. Ademas sabemos que $\widetilde{\varphi}$ es difeomorfismo (es la carta que usamos para probar que era $\Lambda^{k}( T^{*}M)$ era estructura) 6. Luego si $w$ es suave entonces $\widetilde{\varphi}\circ w$ es suave, por lo tanto cada coordenada es suave luego $C_{i_{1}\ldots i_{k}}$ con $1<i_{1}<\ldots<i_{k}\leq n$ son suaves 7. Y obviamente si $C_{i_{1}\ldots i_{k}}$ con $1<i_{1}<\ldots<i_{k}\leq n$ son suaves, $\widetilde{\varphi}\circ w$ es suave. Entonces $w$ lo es (por que si esto es suave componer a derecha con $\varphi ^{-1}:\varphi(U)\rightarrow U$ sigue siendo suave, por que $\varphi ^{-1}$ es homeo, osea que al componer no cambiamos codominio)

\bookqed
\end{bookproof}
\begin{bookdefinition}{Producto cuña de $k$-formas}
Sean $\omega\in\Omega^k(M)$ y $\theta\in\Omega^\ell(M)$. Se define el producto cuña de $\omega$ y $\theta$, denotado por $\omega\wedge\theta$, de manera puntual:

\[
(\omega\wedge\theta)_p:=\omega_p\wedge\theta_p.
\]
En el caso de $f\in C^\infty(M)=\Omega^0(M)$ y $\omega\in\Omega^k(M)$, tenemos

\[
f\wedge\omega=f\omega
\]
osea $(f\wedge w)_{p}=f(p)w_{p}$

\end{bookdefinition}
\begin{bookproposition}{Ejercicio del practico}
Probar que si $\omega\in\Omega^k(M)$ y $\eta\in\Omega^\ell(M)$, entonces $\omega\wedge\eta\in\Omega^{k+\ell}(M)$.

\end{bookproposition}
\begin{bookproof}{}
\begin{enumerate}
\item En coordenadas, escribimos
\end{enumerate}
\[
\omega=\sum_I a_I\,dx^I,\qquad \eta=\sum_J b_J\,dx^J,
\]
con $a_I,b_J$ funciones suaves. 2. Entonces

\[
\omega\wedge\eta=\sum_{I,J}a_Ib_J\,dx^I\wedge dx^J.
\]
\begin{enumerate}[start=3]
\item Los términos que tienen diferenciales repetidos dan $0$, y los demás se reordenan posiblemente cambiando el signo. Pero los nuevos coeficientes siguen siendo productos $\pm a_Ib_J$ de funciones suaves.
\item Por lo tanto, en cada carta $\omega\wedge\eta$ tiene coordenadas suaves. Luego $\omega\wedge\eta\in\Omega^{k+\ell}(M)$.
\end{enumerate}
\bookqed
\end{bookproof}
\phantomsection\label{blk:ce25ad}
\begin{bookdefinition}{Construccion de $\Omega^{*}(M)$}
$\Omega^{k}(M)$ es espacio vectorial sobre $\mathbb{R}$ La suma y producto escalar se definen de manera puntual, $\omega,\theta\in\Omega^k(M)$ y $\lambda\in\mathbb R$,

\[
(\omega+\lambda\theta)_p=\omega_p+\lambda\theta_p,
\]
y se ve que $\omega+\lambda\theta\in\Omega^k(M)$.

Ademas $\Omega^k(M)$ es un $C^\infty(M)$-modulo con la accion dada por: si $f\in C^\infty(M)$ y $\omega\in\Omega^k(M)$, entonces

\[
(f\omega)_p=f(p)\omega_p,
\]
y se ve que $f\omega\in\Omega^k(M)$. Extendiendo el producto cuña por bilinealidad a todo $\Omega^*(M)$, tenemos que este espacio es un espacio vectorial real con estructura de algebra dada por $\wedge$, que es graduada osea tiene una descomposicion en subespacios, en este caso $\Omega^k(M)$, $k=0,\ldots,n$, que cumplen

\[
\Omega^k(M)\wedge\Omega^\ell(M)\subset\Omega^{k+\ell}(M),
\]
y tambien es un $C^\infty(M)$-modulo.

\end{bookdefinition}
\phantomsection\label{blk:76ed6d}
\section{La diferencial exterior}
\begin{bookremark}{Motivacion}
Notar que, dado $f\in C^\infty(M)=\Omega^0(M)$, tenemos un operador $d$, que aplicado a $f$ nos da $df\in\Omega^1(M)$. Queremos extender este operador al resto de espacios de $k$-formas suaves. Que $df\in \Omega^{1}(M)$ es por que $(df)_{p}$ me da un $1$-tensor. Por que $(df)_{p}:T_{p}M\rightarrow T_{f(p)}(\mathbb{R})\simeq \mathbb{R}$ (usando la identificacion con las coordenadas en la imagen) osea $(df)_{p}\in T^{1}(T_{p}^{*}M)=\Lambda^{1}(T_{p}^{*}M)$

\end{bookremark}
\phantomsection\label{blk:171d2a}
\begin{booklemma}{La diferencial exterior es local, lema previo}
Sea $d:\Omega^*(M)\to\Omega^*(M)$ satisfaciendo (i)-(iv). Sean $\omega_1,\omega_2\in\Omega^*(M)$ tales que coinciden en un abierto $U$ de $M$. Entonces

\[
(d\omega_1)|_U=(d\omega_2)|_U.
\]
\end{booklemma}
\begin{bookproof}{}
\begin{enumerate}
\item Como $d$ es transformacion lineal, podemos mirar $\omega=\omega_1-\omega_2$ y asi es suficiente con ver que si $\omega|_U=0$ (Osea $\omega_{1}$ y $\omega_{2}$ coinciden en $U$ abierto), entonces $(d\omega)|_U=0$. Con lo cual se cumpliria la igualdad pedida
\item Aca hacemos un truco clasico de tensores y es volver todo cero usando una funcion de levantamiento.
\item Sabemos que las funciones de levantamiento preservan en un abierto contenido en otro y matan por "afuera", pero aca $\omega$ ya es cero en $U$. La idea es "volverlo cero" en todo $M$.
\item Sea $p\in U$ arbitrario. Sea $V$ abierto de $M$ con $p\in V$ y $\overline V\subset U$. Sea $f\in C^\infty(M)$ funcion de levantamiento tal que $f|_{\overline V}\equiv1$ y $\operatorname{supp}f\subset U$.
\item Como $f$ es suave, entonces $f\omega\in\Omega^*(M)$, ademas
\end{enumerate}
\[
(f\omega)_q=f(q)\omega_q=\begin{cases}f(q)0,&q\in U,\\0\omega_q,&q\notin U,\end{cases}
\]
porque $\operatorname{supp}f\subset U$. Asi $f\omega=0$ en todo $M$. 6. Como $d$ es lineal, $d(f\omega)=d(0)=0$. Pero por otro lado

\[
d(f\omega)=d(f\wedge\omega)=(df)\wedge\omega+(-1)^{0.k} f\,d\omega.
\]
\begin{enumerate}[start=7]
\item Luego
\end{enumerate}
\[
0=(df)\wedge\omega+f\,d\omega.
\]
\begin{enumerate}[start=8]
\item Evaluando en $p\in U$, tenemos
\end{enumerate}
\[
0=((df)\wedge\omega)_p+f(p)(d\omega)_p=(df)_p\wedge\omega_p+(d\omega)_p=(d\omega)_p,
\]
porque $\omega_p=0$ y $f(p)=1$. Como $p$ era arbitrario, $(d\omega)|_U=0$.

\bookqed
\end{bookproof}
\phantomsection\label{blk:ef3cfa}
\begin{booktheorem}{Derivada exterior}
Sea $M$ variedad suave. Entonces existe una unica transformacion lineal

\[
d:\Omega^*(M)\to\Omega^*(M)
\]
llamada derivada exterior, que satisface:

\begin{enumerate}
\item $d:\Omega^k(M)\to\Omega^{k+1}(M)$.
\item $d(f)$ es $df$, diferencial de $f$, para toda $f\in\Omega^0(M)=C^\infty(M)$.
\item Si $\omega\in\Omega^k(M)$ y $\theta\in\Omega(M)$, entonces
\end{enumerate}
\[
d(\omega\wedge\theta)=(d\omega)\wedge\theta+(-1)^k\omega\wedge d\theta.
\]
\begin{enumerate}[start=4]
\item $d^2=0$.
\end{enumerate}
\end{booktheorem}
\begin{bookproof}{}
\begin{itemize}
\item \textbf{Unicidad}
\end{itemize}
\begin{enumerate}
\item La idea de la prueba es construir $d$ de forma local, usando las condiciones que debemos satisfacer: (ii)-(iv). Hacemos ingenieria inversa para saber como deberia ser $d$, localmente. Primero, notemos:
\item Asumiendo existencia veamos primero la unicidad, que nos va dar pistas de como definir $d$.
\item Como $d$ se va definir localmente , fijemos una carta suave $(U,\varphi=(x_1,\ldots,x_n))$ y tomemos $\omega\in\Omega^k(M)$ arbitrario y llegemos a una formula para $(d\omega)|_U$.
\item Tenemos
\end{enumerate}
\[
\omega|_U=\sum_{1\le i_1<\cdots<i_k\le n}C_{i_1\ldots i_k}\,dx_{i_1}\wedge\cdots\wedge dx_{i_k}.
\]
\begin{enumerate}[start=5]
\item Pero $d$ actua sobre cosas definidas sobre todo $M$, entonces no podemos aplicar directamente $d$ a la expresion anterior. Para hacerlo bien, debemos extender dicha expresion a todo $M$, usando una funcion de levantamiento.
\item Sea $V$ abierto de $M$ con $\overline V\subset U$, y sea $f\in C^\infty(M)$ tal que $f|_{\overline V}\equiv1$ y $\operatorname{supp}f\subset U$. Extendemos las $C_{i_1\ldots i_k}$ y $x_1,\ldots,x_n$ definiendo:
\end{enumerate}
\[
f\cdot x_i=\begin{cases}f(q)x_i(q),&q\in U,\\0,&\text{en otro caso},\end{cases}\qquad f\cdot C_{i_1\ldots i_k}=\begin{cases}f(q)C_{i_1\ldots i_k}(q),&q\in U,\\0,&\text{en otro caso}.\end{cases}
\]
\begin{enumerate}[start=7]
\item Estas nuevas funciones son suaves (hemos hecho cuentas similares en otras ocaciones)
\item Consideremos la nueva $k$-forma:
\end{enumerate}
\[
\widetilde{\omega}=\sum_{1\le i_1<\cdots<i_k\le n}\underbrace{f.C_{i_1\cdots i_k}}_{\text{smooth function}}\ \ \underbrace{\underbrace{d(f.x_{i_1})}_{\text{smooth }1\text{-form}}\wedge\cdots\wedge\underbrace{d(f.x_{i_k})}_{\text{smooth }1\text{-form}}}_{\text{smooth }k\text{-form}}.
\]
$d(f.x_{i_{j}})$ son $1$-formas por lo mismo que \hyperref[blk:171d2a]{Motivacion} y luego usamos \hyperref[blk:ce25ad]{Ejercicio del practico} para el wedge de todas esas $1$-formas y llegamos a la $k$-forma 9. Notar que $\widetilde\omega$ coincide con $\omega$ en el abierto $V$. En efecto, si $q\in V$, tenemos $f(q)=1$ y asi

\[
(f\cdot C_{i_1\ldots i_k})(q)=f(q).C_{i_1\ldots i_k}(q)=C_{i_1\ldots i_k}(q)
\]
\begin{enumerate}[start=10]
\item Y ademas $d(f\cdot x_i)_q=(dx_i)_q$ por que
\end{enumerate}
\[
d(f.x_{i})_{q}\left(\frac{\partial}{\partial x_{i}}\bigg|_{q}\right)=\left(\frac{\partial}{\partial x_{i}}\bigg|_{q}\right)(f.x_{i})=\left(\frac{\partial}{\partial x_{i}}\bigg|_{q}\right)(x_{i})=(dx_{i})_{q}\left(\frac{\partial}{\partial x_{i}}\bigg|_{q}\right)
\]
usando definicion de diferencial y en el segundo igual como $f.x_{i}\equiv x_{i}$ en el abierto $V$ sabemos que $\left(\frac{\partial}{\partial x_{i}}\bigg|_{q}\right)$ actua de la misma forma en ambos (germen local) 11. Pero entonces

\[
(d(f.x_{i_{1}})\wedge\ldots\wedge d(f.x_{i_{k}}))_{q}=(d(x_{i_{1}})\wedge\ldots\wedge d(x_{i_{k}}))_{q}
\]
usando la definicion de producto cuña de $k$-formas y la asociatividad 12. Luego $w$ y $\widetilde{w}$  coinciden en un abierto $V$. Por tanto, por \hyperref[blk:ef3cfa]{La diferencial exterior es local, lema previo}

\[
(d\omega)|_V=(d\widetilde\omega)|_V.
\]
\begin{enumerate}[start=13]
\item Lo bueno es que la expresion para $\widetilde\omega$ esta definida sobre todo $M$. Por lo tanto podemos aplicar $d$ y obtenemos:
\end{enumerate}
\[
(d\widetilde\omega)=d\sum fC_{i_1\ldots i_k}\,d(fx_{i_1})\wedge\cdots\wedge d(fx_{i_k})=d\sum fC_{i_1\ldots i_k}\wedge d(fx_{i_1})\wedge\cdots\wedge d(fx_{i_k})
\]
el ultimo igual por que $fC_{i_1\ldots i_k}$ es una $0$-forma 14. Por linealidad de $d$ y usando la propiedad $3$ de esta misma afirmacion y considerando que $fC_{i_1\ldots i_k}$ es una $0$-forma, queda

\[
(d\widetilde\omega)=\sum d(fC_{i_1\ldots i_k})\wedge d(fx_{i_1})\wedge\cdots\wedge d(fx_{i_k})+\sum fC_{i_1\ldots i_k}\wedge d( d(fx_{i_1})\wedge\cdots\wedge d(fx_{i_k}))
\]
\begin{enumerate}[start=15]
\item Ahora como $d^{2}=0$
\end{enumerate}
\[
(d\widetilde\omega)=\sum d(fC_{i_1\ldots i_k})\wedge d(fx_{i_1})\wedge\cdots\wedge d(fx_{i_k})
\]
\begin{enumerate}[start=16]
\item Finalmente como por paso 12. $(d\omega)|_{V}=(d \widetilde{\omega})|_{V}$ y por paso 16. y por la misma idea que en 11. tenemos
\end{enumerate}
\[
(d\omega)|_V=(d \widetilde{\omega})|_{V}=\sum d(C_{i_1\ldots i_k})\wedge dx_{i_1}\wedge\cdots\wedge dx_{i_k}.
\]
\begin{enumerate}[start=17]
\item De lo que resulta que $d$ es unico pues la formula anterior solo depende de la diferencial usual de funciones suaves.
\item Notar que la formula anterior da pistas de como se debe definir la transformacion lineal $d:\Omega^*(M)\to\Omega^*(M)$.
\end{enumerate}
\begin{itemize}
\item \textbf{Existencia}
\end{itemize}
\begin{enumerate}
\item Construyamos $d$ localmente y luego lo hacemos de forma global. Fijemos una carta $(U,\varphi=(x_1,\ldots,x_n))$.
\item Dada una $k$-forma $\omega\in\Omega^k(U)$,
\end{enumerate}
\[
\omega=\sum C_{i_1\ldots i_k}\,dx_{i_1}\wedge\cdots\wedge dx_{i_k},\qquad C_{i_1\ldots i_k}\in C^\infty(U),
\]
\begin{enumerate}[start=3]
\item Definimos $d_U$ por
\end{enumerate}
\[
d_U\omega:=\sum dC_{i_1\ldots i_k}\wedge dx_{i_1}\wedge\cdots\wedge dx_{i_k},
\]
donde $dC_{i_1\ldots i_k}$ es la diferencial de $C_{i_1\ldots i_k}$, por lo tanto una $1$-forma y luego su base es $\{ dx_{1},\ldots,dx_{n} \}$

\[
dC_{i_1\ldots i_k}=\sum_{j=1}^n\frac{\partial C_{i_1\ldots i_k}}{\partial x_j}\,dx_j.
\]
los coeficientes salen por que $dC_{I}\frac{\partial}{\partial x_{j}}=\frac{\partial}{\partial x_{j}}C_{I}$ y esto vale por \textasciicircum{}311f86 4. Por linealidad del diferencial de funciones suaves, claramente $d_U$ se extiende a una transformacion lineal sobre todo $\Omega^*(U)$ 5. Ademas $d_U$ satisface \textbf{(i)} por que si reemplazamos la expresion que dimos en 3. llegamos a

\[
d_U\omega:=\sum_{j=1}^{n}  \sum_{i_{1}< \ldots<i_{k}}\frac{\partial C_{i_1\ldots i_k}}{\partial x_j}\,dx_j\wedge dx_{i_1}\wedge\cdots\wedge dx_{i_k},
\]
y luego de permutar y posiblemente agregar un $-1$ nos queda que esto es una suma de $1$ hasta $n$ de $k+1$-formas 6. \textbf{(ii)} es trivial por que $f=f.1$ entonces $d|_{U}f=d|_{U}f.1=d|_{U}f$ 7. \textbf{(iii)} Como $d_U$ es lineal y $\wedge$ es bilineal, basta con analizar cuando $\omega=f\,dx_{i_1}\wedge\cdots\wedge dx_{i_k}$ y $\theta=g\,dx_{j_1}\wedge\cdots\wedge dx_{j_\ell}$. Esto es por que el producto cuña de dos elementos de la suma directa se define como la suma de los productos cuña de los pares son $k$-formas del mismo grado $k$. 8. Luego por linealidad seria diferencial de cada uno de esos productos. Pero esos productos son basicamente como en multiplicacion de polinomios la suma de todas las combinaciones, y luego por linealidad nuevamente basta verlo para un caso cualquiera de esos que seria $\omega\wedge\theta$ 9. Entonces

\[
\begin{aligned}d_U(\omega\wedge\theta)& =d_U(fg\,dx_{i_1}\wedge\cdots\wedge dx_{i_k}\wedge dx_{j_1}\wedge\cdots\wedge dx_{j_\ell})\\ & =d(fg)\wedge dx_{i_1}\wedge\cdots\wedge dx_{i_k}\wedge dx_{j_1}\wedge\cdots\wedge dx_{j_\ell}\\&=(f\,dg+g\,df)\wedge dx_{i_1}\wedge\cdots\wedge dx_{i_k}\wedge dx_{j_1}\wedge\cdots\wedge dx_{j_\ell}\\& =f\,dg\wedge dx_{i_1}\wedge\cdots\wedge dx_{i_k}\wedge dx_{j_1}\wedge\cdots\wedge dx_{j_\ell}+g\,df\wedge dx_{i_1}\wedge\cdots\wedge dx_{i_k}\wedge dx_{j_1}\wedge\cdots\wedge dx_{j_\ell}\\&=(-1)^k f\,dx_{i_1}\wedge\cdots\wedge dx_{i_k}\wedge dg\wedge dx_{j_1}\wedge\cdots\wedge dx_{j_\ell}+df\wedge dx_{i_1}\wedge\cdots\wedge dx_{i_k}\wedge g\,dx_{j_1}\wedge\cdots\wedge dx_{j_\ell}\\&=(d_U\omega)\wedge\theta+(-1)^k\omega\wedge d_U\theta.\end{aligned}
\]
\begin{enumerate}[start=10]
\item \textbf{(iv)} $d_U^2=0$. Por linealidad de $d_U$, es suficiente con ver $d_Ud_U\omega=0$ con $\omega=f\,dx_I$, donde $dx_I=dx_{i_1}\wedge\cdots\wedge dx_{i_k}$.
\end{enumerate}
\[
\begin{aligned}d_Ud_U\omega&=d_U\left(\sum_j\frac{\partial f}{\partial x_j}dx_j\wedge dx_I\right)\\&=\sum_j d\left(\frac{\partial f}{\partial x_j}\right)\wedge dx_j\wedge dx_I\\&=\sum_j\sum_k\frac{\partial^2 f}{\partial x_k\partial x_j}dx_k\wedge dx_j\wedge dx_I\\&=\sum_{1\le j<k\le n}\left(\frac{\partial^2 f}{\partial x_k\partial x_j}-\frac{\partial^2 f}{\partial x_j\partial x_k}\right)dx_k\wedge dx_j\wedge dx_I=0.\end{aligned}
\]
Esto último vale porque

\[
\left(\frac{\partial^2 f}{\partial x_k\partial x_j}-\frac{\partial^2 f}{\partial x_j\partial x_k}\right)=\left[\frac{\partial }{\partial x_k},\frac{\partial}{\partial x_j}\right]f=0.
\]
\begin{enumerate}[start=11]
\item Por ultimo, definimos de manera global $d$. Sea $\theta\in\Omega^r(M)$. Dado $p\in M$, definimos
\end{enumerate}
\[
(d\theta)_p:=(d_U\theta|_U)_p
\]
para cualquier carta suave $(U,\varphi=(x_1,\ldots,x_n))$ de $M$ alrededor de $p$. (la parte. dela derecha cae en lo que hicimos al principio por lo tanto cumple todo) 12. Necesitamos ver que esta bien definido. Sea $(V,\psi=(y_1,\ldots,y_n))$ otra carta suave de $p$, y $W=U\cap V\ne\varnothing$. Por linealidad de la derivada exterior, alcanza con verlo cuando $\theta\in\Omega^k(M)$. 13. Tenemos

\[
\theta|_U=\sum_Ia_I\,dx_I,\qquad \theta|_V=\sum_Jb_J\,dy_J.
\]
\begin{enumerate}[start=14]
\item Entonces por definicion de $d_U$ y $d_V$,
\end{enumerate}
\[
d_U\theta|_U=\sum_Ida_I\wedge dx_I,\qquad d_V\theta|_V=\sum_Jdb_J\wedge dy_J.
\]
\begin{enumerate}[start=15]
\item Pero en $W=U\cap V$ hay una unica diferencial exterior local
\end{enumerate}
\[
d_W:\Omega^*(W)\to\Omega^{*}(W).
\]
\begin{enumerate}[start=16]
\item Tomando las dos expresiones de $\theta|_W$, esto ultimo es igual a, por ser $d_W$ diferencial exterior,
\end{enumerate}
\[
\sum_Ida_I\wedge dx_I=\sum_Jdb_J\wedge dy_J.
\]
\begin{enumerate}[start=17]
\item En particular, si evaluamos en $p$,
\end{enumerate}
\[
(\sum_Ida_I\wedge dx_I)_p=(\sum_Jdb_J\wedge dy_J)_p.
\]
\bookqed
\end{bookproof}
\phantomsection\label{blk:f7b034}
\section{Glosario}
\begin{bookremark}{No se de donde salio esto}
Si $\omega\in\Lambda^kT_p^*M$, entonces $\omega$ es una aplicación $k$-lineal alternada

\[
\omega:T_pM\times\cdots\times T_pM\to\mathbb R.
\]
La carta $(U,\varphi=(x_1,\ldots,x_n))$ induce una base de $T_pM$ dada por

\[
\left\{\left.\frac{\partial}{\partial x_1}\right|_p,\ldots,\left.\frac{\partial}{\partial x_n}\right|_p\right\}.
\]
Por eso, para cada $1\le i_1<\cdots<i_k\le n$, se define

\[
v_{i_1\cdots i_k}(\omega):=\omega\left(\left.\frac{\partial}{\partial x_{i_1}}\right|_p,\ldots,\left.\frac{\partial}{\partial x_{i_k}}\right|_p\right).
\]
Es decir, $v_{i_1\cdots i_k}(\omega)$ es el coeficiente de $\omega$ obtenido al evaluarla en los vectores coordenados correspondientes.

\end{bookremark}
\phantomsection\label{blk:4e8085}
